\chapter{Manual de instalare și utilizare}
\pagestyle{fancy}

Acest capitol descrie pașii necesari pentru instalarea și utilizarea aplicației PyVisionAnnotator. Aplicația poate fi folosită în două moduri: ca program desktop împachetat în executabil Windows sau ca proiect rulat din surse. Modul local permite încărcarea imaginilor, crearea adnotărilor și exportul datasetului, iar modul colaborativ necesită pornirea backend-ului Docker pentru autentificare, chat și sincronizare în timp real.

\section{Cerințe hardware}

Cerințele hardware depind de scenariul de utilizare. Adnotarea manuală are cerințe reduse, dar fluxurile AutoSeg, AutoMask și chatbotul local pot consuma mai multă memorie și pot beneficia de GPU.

\begin{table}[H]
    \centering
    \small
    \begin{tabular}{p{0.26\textwidth} p{0.30\textwidth} p{0.34\textwidth}}
        \hline
        Componentă & Minim & Recomandat \\
        \hline
        Client desktop & Procesor dual-core, 4 GB RAM (Random Access Memory) & Procesor quad-core, 8 GB RAM sau mai mult \\
        Spațiu pe disc & 2 GB pentru aplicație și imagini de test & 10 GB sau mai mult pentru dataseturi și exporturi \\
        Backend Docker & 4 GB RAM disponibili & 8 GB RAM disponibili pentru rulare stabilă cu toate serviciile \\
        AutoSeg/AutoMask & CPU modern & GPU NVIDIA cu drivere actualizate, dacă modelele externe folosesc CUDA (Compute Unified Device Architecture) \\
        LLM local & 8 GB RAM pentru modele mici & 16 GB RAM sau mai mult pentru modele locale mai mari \\
        \hline
    \end{tabular}
    \caption{Cerințe hardware minime și recomandate}
    \label{tab:cerinte-hardware}
\end{table}

\section{Cerințe software}

Pentru utilizatorul obișnuit, varianta recomandată este rularea aplicației din executabilul generat. În acest caz, nu este necesară instalarea separată a pachetelor Python. Pentru dezvoltare, testare sau reconstruirea aplicației sunt necesare Python și dependențele proiectului.

\begin{table}[H]
    \centering
    \small
    \begin{tabular}{p{0.30\textwidth} p{0.58\textwidth}}
        \hline
        Scenariu & Software necesar \\
        \hline
        Rulare din executabil & Windows 64-bit și folderul generat \texttt{build\_output/PyVisionAnnotator}. \\
        Rulare din surse & Python 3, pip, mediul virtual și pachetele \texttt{PyQt6}, \texttt{websocket-client}, \texttt{ollama}, \texttt{llama-cpp-python}. \\
        Build aplicație desktop & PyInstaller, instalat automat de \texttt{utils/build\_exe.py} dacă lipsește. \\
        Build modele externe & PyInstaller și scriptul \texttt{research/others/exporter.py}. \\
        Backend colaborativ & Docker Desktop, Docker Compose și fișierul \texttt{.env} din \texttt{ColaborativeServer}. \\
        Chatbot local & Ollama instalat și pornit, dacă se folosește backend-ul Ollama. \\
        \hline
    \end{tabular}
    \caption{Cerințe software pe scenarii de rulare}
    \label{tab:cerinte-software}
\end{table}

\section{Instalarea aplicației din executabil}

Pentru distribuirea aplicației desktop se folosește scriptul \texttt{AnnotationEngine/utils/build\_exe.py}. Acesta pornește PyInstaller, construiește aplicația în modul \texttt{onedir} și creează un pachet complet în directorul \texttt{AnnotationEngine/build\_output}. La final, executabilul principal se găsește în subdirectorul \texttt{PyVisionAnnotator}.

\begin{verbatim}
cd AnnotationEngine
python utils/build_exe.py
\end{verbatim}

După rularea comenzii, aplicația se pornește prin deschiderea fișierului:

\begin{verbatim}
AnnotationEngine/build_output/PyVisionAnnotator/PyVisionAnnotator.exe
\end{verbatim}

Scriptul de build copiază și fișierele auxiliare disponibile, precum \texttt{requirements.txt}. Dacă directorul \texttt{research/others/dist} există, acesta este copiat în aplicația împachetată sub numele \texttt{inference\_backends}. La prima pornire, aplicația caută automat în acest folder executabilele pentru AutoMask și AutoSeg, precum și fișierele de weights cu extensia \texttt{.pt}. Dacă aceste fișiere sunt găsite, căile sunt salvate automat în setările aplicației.

\section{Împachetarea modelelor externe}

Fluxurile AutoSeg și AutoMask pot rula scripturi Python sau executabile. Pentru a distribui mai ușor modelele externe, în directorul \texttt{research/others} există scriptul \texttt{exporter.py}. Acesta folosește PyInstaller cu opțiunile \texttt{--onedir}, \texttt{--clean} și \texttt{--noupx}, iar rezultatul este creat în folderul \texttt{dist}.

Înainte de rulare, variabila \texttt{TARGET\_SCRIPT} din \texttt{exporter.py} trebuie setată la scriptul dorit, de exemplu \texttt{mask2former.py} sau \texttt{yoloe.py}. Apoi se rulează:

\begin{verbatim}
cd research/others
python exporter.py
\end{verbatim}

Executabilul generat poate fi folosit direct din setările aplicației. Pentru un pachet complet, folderul \texttt{dist} rezultat trebuie să existe înainte de rularea \texttt{utils/build\_exe.py}, astfel încât builder-ul aplicației să îl copieze automat în \texttt{inference\_backends}. Dacă se preferă configurarea manuală, utilizatorul poate deschide dialogul de setări și poate selecta căile către executabilul AutoMask, executabilul AutoSeg și weights-ul YOLO.

\section{Instalarea din surse}

Rularea din surse este utilă pentru dezvoltare sau pentru verificarea rapidă a aplicației fără a crea executabilul. Pașii de bază sunt:

\begin{verbatim}
cd AnnotationEngine
python -m venv .venv
.venv\Scripts\activate
python -m pip install -r requirements.txt
python app.py
\end{verbatim}

Aplicația pornește fereastra principală PyVisionAnnotator. Dacă se folosesc AutoSeg, AutoMask sau chatbot local, trebuie instalate și dependințele necesare modelelor respective în mediul în care rulează scripturile externe. Pentru Ollama, aplicația are nevoie de pachetul Python \texttt{ollama}, dar și de serviciul Ollama pornit separat pe sistem.

\section{Pornirea backend-ului colaborativ}

Adnotarea locală și exportul fișierelor pot fi folosite fără backend. Pentru autentificare, chat și colaborare în timp real este necesar serverul din \texttt{ColaborativeServer}. Acesta se pornește prin Docker Compose:

\begin{verbatim}
cd ColaborativeServer
docker compose up --build
\end{verbatim}

Serverul pornește Traefik, RabbitMQ, Postgres și cele trei servicii Spring Boot: autentificare, chat și annotation. Fișierul \texttt{.env} trebuie să existe în directorul \texttt{ColaborativeServer} și să conțină valorile pentru utilizatorii și parolele serviciilor, portul serviciului de autentificare și secretul JWT.

După pornire, clientul folosește implicit adresa \texttt{http://localhost}. Rutele principale sunt disponibile prin proxy: \texttt{/auth} pentru cont, \texttt{/ws} pentru chat și \texttt{/annotation} pentru colaborare. Dacă se schimbă hostul sau portul, adresele pot fi modificate din dialogul de setări al aplicației.

\begin{table}[H]
    \centering
    \small
    \begin{tabular}{p{0.32\textwidth} p{0.56\textwidth}}
        \hline
        Comandă & Rol \\
        \hline
        \texttt{python utils/build\_exe.py} & Construiește aplicația desktop în \texttt{build\_output}. \\
        \texttt{python exporter.py} & Construiește executabilul pentru scriptul extern ales în \texttt{TARGET\_SCRIPT}. \\
        \texttt{python app.py} & Pornește aplicația din surse. \\
        \texttt{docker compose up --build} & Pornește backend-ul colaborativ. \\
        \hline
    \end{tabular}
    \caption{Comenzi de instalare, build și pornire}
    \label{tab:comenzi-instalare}
\end{table}

\section{Utilizarea aplicației}

La pornire, aplicația deschide fereastra principală, formată din lista de imagini, canvasul central și panoul de unelte/proprietăți. Utilizatorul alege un folder, selectează imaginea dorită și poate naviga prin zoom cu rotița mouse-ului sau prin deplasare cu click mijlociu ori Space și click stânga.

Adnotările se creează din panoul de unelte: dreptunghi pentru bounding box, poligon prin puncte succesive și mască prin AutoMask, brush sau eraser. Elementul selectat poate fi redenumit, recolorat sau șters din panoul de proprietăți. Salvarea internă se face în JSON, iar exporturile disponibile includ formatele COCO și YOLO. Dacă există modificări nesalvate la schimbarea imaginii sau la închiderea aplicației, utilizatorul primește un dialog de confirmare.

AutoSeg, AutoMask și endpoint-urile de cont se configurează din fereastra de setări. Pentru AutoSeg se selectează executabilul YOLO și fișierul de weights \texttt{.pt}, iar pentru AutoMask se selectează executabilul sau scriptul de segmentare pe punct. Pentru colaborare, utilizatorul se autentifică din butonul de cont, deschide fereastra de chat și creează sau accesează o sesiune comună. Chatbotul local poate fi folosit separat prin Ollama sau llama.cpp, cu condiția ca modelul local ales să fie disponibil.

\section{Probleme frecvente}

Dacă autentificarea, chatul sau colaborarea nu funcționează, se verifică mai întâi backend-ul Docker și adresa configurată în aplicație, implicit \texttt{http://localhost}. Pentru probleme AutoSeg sau AutoMask, se verifică existența executabilului/scriptului, a modelului configurat și, în pachetul final, a folderului \texttt{inference\_backends}. Pentru rularea din surse, cele mai frecvente cauze sunt mediul virtual neactivat, pachetele din \texttt{requirements.txt} neinstalate sau serviciul Ollama nepornit.
